\documentclass[
	% -- opções da classe memoir --
	12pt,				% tamanho da fonte
	%openright,			% capítulos começam em pág ímpar (insere página vazia caso preciso)
	oneside,   	        % para impressão em verso e anverso use twoside. Oposto a oneside
	a4paper,			% tamanho do papel. 
	% -- opções da classe abntex2 --
	%chapter=TITLE,		% títulos de capítulos convertidos em letras maiúsculas
	%section=TITLE,		% títulos de seções convertidos em letras maiúsculas
	%subsection=TITLE,	% títulos de subseções convertidos em letras maiúsculas
	%subsubsection=TITLE,% títulos de subsubseções convertidos em letras maiúsculas
	% -- opções do pacote babel --
	english,			% idioma adicional para hifenização
	french,				% idioma adicional para hifenização
	spanish,			% idioma adicional para hifenização
	brazilian,				% o último idioma é o principal do documento
	]{pacotes/abntex2}


% ---
% PACOTES
% ---

% ---
% Pacotes fundamentais 
% ---
\usepackage{cmap}				% Mapear caracteres especiais no PDF
\usepackage{lmodern}			% Usa a fonte Latin Modern
\usepackage[T1]{fontenc}		% Selecao de codigos de fonte.
\usepackage[utf8]{inputenc}		% Codificacao do documento (conversão automática dos acentos)
\usepackage{indentfirst}		% Indenta o primeiro parágrafo de cada seção.
\usepackage{color}				% Controle das cores
\usepackage{graphicx}			% Inclusão de gráficos
\setlength{\arrayrulewidth}{0.5mm}
\setlength{\tabcolsep}{18pt}
\renewcommand{\arraystretch}{1.5}
% ---

% ---
% Pacotes adicionais, usados no anexo do modelo de folha de identificação
% ---
\usepackage{multicol}
\usepackage{multirow}
\usepackage[table]{xcolor}
\usepackage{longtable}
% ---
	
% ---
% Pacotes adicionais, usados apenas no âmbito do Modelo Canônico do abnteX2
% ---
\usepackage{lipsum}				% para geração de dummy text
% ---

% ---
% Pacotes de citações
% ---
\usepackage[brazilian,hyperpageref]{backref}	 % Paginas com as citações na bibl
\usepackage[alf]{pacotes/abntex2cite}	% Citações padrão ABNT
\usepackage{comment}

% Estilização da linguagem PLY
\usepackage{listings}
\usepackage{xcolor}

% VS Code Light+ colors
\definecolor{vscKeyword}{HTML}{0000FF}
\definecolor{vscComment}{HTML}{008000}
\definecolor{vscString}{HTML}{A31515}
\definecolor{vscFunction}{HTML}{795E26}
\definecolor{vscNumbers}{HTML}{808080}
\definecolor{vscBackground}{HTML}{FFFFFF}

\lstdefinestyle{ply}{
    aboveskip=0.8\baselineskip,
    backgroundcolor=\color{vscBackground},
    basicstyle=\ttfamily\small,
    numbers=left,
    numberstyle=\tiny\color{vscNumbers},
    numbersep=8pt,
    breaklines=true,
    keepspaces=true,
    showspaces=false,
    showstringspaces=false,
    showtabs=false,
    tabsize=4,
    frame=single,
    rulecolor=\color{black!20},
    keywordstyle=\color{vscKeyword}\bfseries,
    commentstyle=\color{vscComment}\itshape,
    stringstyle=\color{vscString},
    captionpos=b
}

\lstset{style=ply}


% --- 
% CONFIGURAÇÕES DE PACOTES
% --- 

% ---
% Configurações do pacote backref
% Usado sem a opção hyperpageref de backref
\renewcommand{\backrefpagesname}{Citado na(s) página(s):~}
% Texto padrão antes do número das páginas
\renewcommand{\backref}{}
% Define os textos da citação
\renewcommand*{\backrefalt}[4]{
	\ifcase #1 %
		Nenhuma citação no texto.%
	\or
		Citado na página #2.%
	\else
		Citado #1 vezes nas páginas #2.%
	\fi}%
% ---

% ---
% Informações de dados para CAPA e FOLHA DE ROSTO
% ---
\titulo{Atividade: Relatório com as Expressões Regulares para a Análise Léxica do seu Projeto}
\autor{Maria Eduarda Bambini}
\local{Campo Mourão}
\data{Setembro / 2026}
\instituicao{%
  Implementação de Linguagens de Programação - OPT011
  \par
  Prof. Rogério Aparecido Gonçalves}
\tipotrabalho{Relatório técnico}
% O preambulo deve conter o tipo do trabalho, o objetivo, 
% o nome da instituição e a área de concentração
% ---

% ---
% Configurações de aparência do PDF final

% alterando o aspecto da cor azul
\definecolor{blue}{RGB}{41,5,195}

% informações do PDF
\makeatletter
\hypersetup{
     	%pagebackref=true,
		pdftitle={\@title}, 
		pdfauthor={\@author},
    	pdfsubject={\imprimirpreambulo},
	    pdfcreator={LaTeX with abnTeX2},
		pdfkeywords={abnt}{latex}{abntex}{abntex2}{relatório técnico}, 
		colorlinks=true,       		% false: boxed links; true: colored links
    	linkcolor=blue,          	% color of internal links
    	citecolor=blue,        		% color of links to bibliography
    	filecolor=magenta,      		% color of file links
		urlcolor=blue,
		bookmarksdepth=4
}
\makeatother
% --- 

% --- 
% Espaçamentos entre linhas e parágrafos 
% --- 

% O tamanho do parágrafo é dado por:
\setlength{\parindent}{1.3cm}

% Controle do espaçamento entre um parágrafo e outro:
\setlength{\parskip}{0.2cm}  % tente também \onelineskip

% ---
% compila o indice
% ---
\makeindex
% ---

% Omite a numeração de capítulos
\renewcommand*\thesection{\arabic{section}}

% ----
% Início do documento
% ----
\begin{document}

% Retira espaço extra obsoleto entre as frases.
\frenchspacing 

% Folha de rosto
% (o * indica que haverá a ficha bibliográfica)
\imprimirfolhaderosto

% ---
% inserir o sumario
% ---
\pdfbookmark[0]{\contentsname}{toc}
\tableofcontents*
\cleardoublepage
% ---

% ELEMENTOS TEXTUAIS
\textual

\makeatletter
\renewcommand{\chapter}{\@gobbletwo}
\makeatother

\section{Introdução}
\label{sec:introducao}
    
Para realizar o mapeamento de expressões regulares para a identificação das classes de \textit{tokens} da linguagem TPP, foram realizadas algumas mudanças na tabela original \cite{regex-tokens} proposta como base. Verifique na tabela \ref{tab:tokens_regex} a escolha final.

\begin{longtable}{|>{\centering\arraybackslash}p{6cm}|>{\centering\arraybackslash}p{6cm}|}
\caption{Tokens e suas respectivas expressões regulares.}
\label{tab:tokens_regex} \\
\hline
\rowcolor{gray!20}
Token & Expressão Regular \\
\hline
\endfirsthead

\hline
\rowcolor{gray!20}
Token & Expressão Regular \\
\hline
\endhead

MAIS & \verb|r'\+'| \\
\hline
MENOS & \verb|r'-'| \\
\hline
VEZES & \verb|r'\*'| \\
\hline
DIVIDE & \verb|r'/'| \\
\hline
DOIS\_PONTOS & \verb|r':'| \\
\hline
VIRGULA & \verb|r','| \\
\hline
MENOR & \verb|r'<'| \\
\hline
MAIOR & \verb|r'>'| \\
\hline
IGUAL & \verb|r'='| \\
\hline
DIFERENTE & \verb|r'<>'| \\
\hline
MENOR\_IGUAL & \texttt{r'<='} \\
\hline
MAIOR\_IGUAL & \texttt{r'>='} \\
\hline
E & \verb|r'&&'| \\
\hline
OU & \texttt{r'||'} \\
\hline
NAO & \verb|r'!'| \\
\hline
ABRE\_PARANTESE & \texttt{r'\textbackslash('} \\
\hline
FECHA\_PARENTESE & \texttt{r'\textbackslash)'} \\
\hline
ABRE\_COLCHETE & \texttt{r'\textbackslash['} \\
\hline
FECHA\_COLCHETE & \texttt{r'\textbackslash]'} \\
\hline
SE & \verb|r'se'| \\
\hline
ENTAO & \verb|r'então'| \\
\hline
FIM & \verb|r'fim'| \\
\hline
REPITA & \verb|r'repita'| \\
\hline
ATE & \verb|r'até'| \\
\hline
ATRIBUICAO & \verb|r':='| \\
\hline
LEIA & \verb|r'leia'| \\
\hline
ESCREVA & \verb|r'escreva'| \\
\hline
RETORNA & \verb|r'retorna'| \\
\hline
INTEIRO & \verb|r'inteiro'| \\
\hline
FLUTUANTE & \verb|r'flutuante'| \\
\hline
NUM\_INTEIRO & \texttt{r'\textbackslash d+'} \\
\hline
NUM\_PONTO\_FLUTUANTE & \texttt{r'\textbackslash d*\textbackslash.\textbackslash d+'} \\
\hline
NUM\_NOTACAO\_CIENTIFICA & \texttt{r'\textbackslash d*\textbackslash.?\textbackslash d+[eE][-+]?\textbackslash d+'} \\
\hline

\end{longtable}

\section{Procedimentos}
\label{(sec:procedimentos}

Note que apenas as expressões regulares referentes ao número de ponto flutuante e número em notação científica sofreram alterações. A expressão regular relativa ao número de ponto flutuante estava definida como \texttt{r'\textbackslash d+\textbackslash.?\textbackslash d*'}, porém isso permite que não exista o símbolo ``.'' separador de parte fracionária do número e também que não exista nenhum número após o símbolo separador (``.''). Além disso, a expressão exige que exista um número antes da parte fracionária. 

A expressão foi alterada para inverter todos esses comportamentos citados, ou seja, não exige que haja um dígito antes da parte fracionária, visto que, por exemplo, o número de ponto flutuante .89 será convertido para 0.89; exige que exista o símbolo ``.'' separador da parte fracionária do número e exige que exista um número após esse símbolo, pois o número 89., por exemplo, estaria incorreto. Dessa forma, a expressão regular após os ajustes é a seguinte: \texttt{r'\textbackslash d*\textbackslash.\textbackslash d+'}, note que o primeiro metacaratere ``d'' agora pode ter zero ou múltiplas correspondências, o símbolo ``.'' agora é obrigatório, e o metacaractere ``d'' referente à parte fracionária precisa ter pelo menos uma correspondência.

A expressão regular referente ao número em notação científica também sofreu modificações. A parte da expressão que diz respeito à mantissa foi alterada de modo a reconhecer números inteiros ou de ponto flutuante, então sofreu alterações similares à expressão regular relacionada ao número de ponto flutuante. Além disso, a parte da expressão regular referente à ordem de grandeza do número, agora reconhece tanto ``e'' quanto ``E'' para indicar o expoente da ordem de grandeza. Essa é a expressão regular após esses ajustes: \texttt{r'\textbackslash d*\textbackslash.?\textbackslash d+[eE][-+]?\textbackslash d+'}, que também pode ser verificado na tabela \ref{tab:tokens_regex}.

% ELEMENTOS PÓS-TEXTUAIS
\postextual
% Referências bibliográficas
\renewcommand{\bibsection}{%
\section{\bibname}
\bibmark
%\ifnobibintoc\else
%\phantomsection
%\addcontentsline{toc}{section}{\bibname}
%\fi
\prebibhook}

\bibliography{references}
\end{document}
